%% ma: L10-B02-0039
%% lop: 10
%% chuong: 1
%% bai: 2
%% dang: 10.2.7
%% loai: TLN
%% muc: VD
%% dap_an: "2"
%% nguon: "Ảnh giáo viên gửi 10/10/2026 (ThS. Nguyễn Hoàng Việt, Chương 2 Tập hợp), Đề 2 (đề thứ hai, không ghi tiêu đề), Phần 3 Câu 17"
%% kien_thuc: "Bài 2 (tập hợp và các phép toán trên tập hợp)"
%% trang_thai: chinh-thuc
%% ghi_chu: "Hiểu A=[m+1;2m-1] là đoạn khác rỗng (m ≥ 2). Nếu cho phép A rỗng (m < 2) thì có vô số giá trị m; đề gốc không nói rõ, đáp số 2 theo cách hiểu đoạn khác rỗng"
\begin{ex}
Cho hai tập hợp $A=[m+1;2m-1]$, $B=(0;6)$. Có bao nhiêu giá trị $m$ nguyên để $A\subset B$.
\shortans{2}
\loigiai{Để $A=[m+1;2m-1]$ là một đoạn cần $m+1\le2m-1\Leftrightarrow m\ge2$. $A\subset B\Leftrightarrow m+1>0$ và $2m-1<6\Leftrightarrow m>-1$ và $m<\dfrac72$. Kết hợp: $2\le m<\dfrac72$, $m$ nguyên: $m\in\{2;3\}$, có $2$ giá trị.}
\end{ex}
